\documentclass[a4paper,landscape]{article}
\usepackage[utf8]{inputenc}
\usepackage[english]{babel}
\usepackage{graphicx}
\usepackage{geometry}
	\geometry{a4paper,
	total={170mm,257mm},
	left=20mm,
	top=20mm,
	right=20mm,
	bottom=20mm,
}
\usepackage[dvipsnames]{xcolor}
\definecolor{gateoxide}{gray}{0.9}
\definecolor{NormalGray}{gray}{0.8}
\definecolor{DarkGray}{gray}{0.7}
\usepackage{tikz}
\usetikzlibrary{arrows,automata,shapes}
\tikzstyle{decision} = [diamond, draw, fill=blue!20, text width=4.5em, text badly centered, node distance=2.5cm, inner sep=0pt]
\tikzstyle{block} = [rectangle, draw, fill=blue!20, text width=5em, text centered, rounded corners, minimum height=4em]
\tikzstyle{line} = [draw, very thick, color=black!50, -latex']
\tikzstyle{cloud} = [draw, ellipse,fill=red!20, node distance=2.5cm,minimum height=2em]

\author{David Lanzendörfer}
\title{Generic LibreSilicon process HKUST (NFF)}
\begin{document}
\maketitle

\begin{abstract}
	Copyright © 2017 LANCEVILLE TECHNOLOGY GROUP CO., LIMITED. All rights reserved. \\

	This process is licensed under the Libre Silicon public license; you can redistribute it and/or modify it under the terms of the Libre Silicon public license
	as published by the Libre Silicon alliance, either version 1 of the License, or (at your option) any later version.

	This design is distributed in the hope that it will be useful, but WITHOUT ANY WARRANTY; without even the implied warranty of MERCHANTABILITY or FITNESS FOR A PARTICULAR PURPOSE.
	See the Libre Silicon Public License for more details. \\

	This is the specification of the free silicon manufacturing standard for manufacturing the ls180nm\footnotemark standard logic cells and related free technology nodes from the LibreSilicon project.
	
	For further clarification consult the complete documentation of the process.

	\footnotetext[1]{https://github.com/leviathanch/ls180nm}
\end{abstract}
\vfill
\newpage

Process Flow of Lanceville Technologies LibreSilicon 180nm

\begin{itemize}
	\item Project: LibreSilicon 180nm
	\item Name: Lanceville Technologies Group
	\item Substrate: P-Substrate silicon wafer <100>
	\item Date: \today
\end{itemize}

\begin{tikzpicture}[node distance = 3cm, auto, thick,scale=0.4, every node/.style={transform shape}]
	% substrate
	\fill[YellowOrange] (0,0) rectangle (15,2);
	\node at (2,0.5) {Si (p-type)};
	% oxide
	\fill[gray] (0,1.6) rectangle (15,2.6);
	\node at (2,2.1) {SiO2};
\end{tikzpicture} \tiny\begin{tabular}{|c|}
\hline
\textbf{Wafer Cleanliness} \\
\hline
Clean \\
\hline
Clean \\
\hline
Clean \\
\hline
Clean \\
\hline
\end{tabular} \begin{tabular}{|c|c|c|c|c|c|}
\hline
\textbf{Step Number} &
\textbf{Equipment} &
\textbf{Location} &
\textbf{Cleanliness} &
\textbf{Process} &
\textbf{Requirements} \\
\hline

0.1 &
A3: Sulfuric Cleaning &
P201000 &
Clean &
Initial Clean &
H2SO4 + H2O2, 10mins, 120C \\

\hline

0.2 &
A2: HF:H2O (1:50) &
P201000 &
Clean &
HF dip &
1 min \\

\hline

0.3 &
Spin Dryer-A &
P201000 &
Clean &
Dry the wafer automatically &
{}
\\

\hline

0.4 &
?? &
?? &
?? &
Silicon dioxide growth &
1000nm \\

\hline

\end{tabular}









\end{document}